\section{N-Capas y atributos de calidad}

\subsection{Enunciado}
La Arquitectura N-Capas es criticada por su bajo rendimiento en características como «Agilidad
general», «Facilidad de despliegue» y «Escalabilidad» (Richards, 2015). Un estudiante propone
que, con el uso de capas abiertas (\emph{open layers}) y técnicas como el «fast-lane reader», se
pueden mitigar estos problemas.

Analice esta afirmación. ¿Cómo afecta la introducción de capas abiertas al principio de
Aislamiento de Capas y a la mantenibilidad a largo plazo? ¿Por qué la solución de capas abiertas
puede llevar a un «sinkhole anti-pattern» y qué dice la regla del 80/20 sobre esto? Finalmente,
explique por qué, para un sistema con requisitos de alta escalabilidad y agilidad, una
arquitectura como la Hexagonal o la Limpia, que se basan en los principios de DDD y
desacoplamiento, son inherentemente más adecuadas que una Arquitectura N-Capas, incluso con
optimizaciones.

\subsection{Respuesta}
\textit{[Redactar con \parencite{richards2015patterns,richards2020fundamentals,martin2017clean}.]}
